\documentclass[a5paper, 10pt]{article}


\usepackage{cite, amsmath, amssymb, booktabs, lastpage, graphicx, float}
\usepackage[hidelinks]{hyperref}

\usepackage{geometry}
\geometry{margin=1.5cm,a5paper,twoside,inner=2cm}

\usepackage{setspace}


\usepackage{amsthm}

\theoremstyle{plain}
\newtheorem{lemma}{Lemma}
\newtheorem{theorem}{Theorem}
\newtheorem{conjecture}{Conjecture}
\newtheorem{proposition}[theorem]{Proposition}
\newtheorem{corollary}[theorem]{Corollary}
\renewcommand{\qedsymbol}{\rule{.7em}{.7em}}

\theoremstyle{remark}
\newtheorem{remark}{Remark}
\newtheorem*{note}{Note}
\newtheorem{case}{Case}

\theoremstyle{definition}
\newtheorem{definition}{Definition}
\newtheorem{example}{Example}

%% Changing the title, author, etc.
\usepackage{graphicx}
\makeatletter
\renewcommand{\maketitle}{
    \begin{center}
    \rule[.2em]{\textwidth}{0.353mm}
        \begin{minipage}[m]{0.35\textwidth}
            {\scriptsize
                \begin{center}
                    \textsf{\textbf{\huge MATCH}\\
                        \textit{Communications in Mathematical\\
                            and in Computer Chemistry}
                    }
            \end{center}}
        \end{minipage}\hfill
        \begin{minipage}[m]{0.65\textwidth}
            \begin{flushright}
                \baselineskip=10pt
                {\scriptsize\sffamily{\itshape  MATCH Commun. Math. Comput. Chem.}
                \textbf{\vol{}}
                (\pubyear)
                    \thepage--\pageref{LastPage}}\\
                {\scriptsize\sffamily {\bfseries ISSN:} 0340--6253}\\
                {\scriptsize\sffamily \textbf{doi:} \doi{10.46793/match.\vol-\num.\id}}
            \end{flushright}
        \end{minipage}
        \rule[1em]{\textwidth}{.353mm}
        \baselineskip=0.30in
        {\Large\bfseries \@title} \par
        \vspace{5mm}
        \baselineskip=0.2in
        {\large\bfseries \@author}\par
        \vspace{1mm}
        {\it \@address} \par
        {\small\tt \@email} \par
        \vspace{3mm}
        {\small (Received \@date)} \par
    \end{center}
    \vspace{3mm}
}
\newcommand{\address}[1]{\def\@address{#1}}
\newcommand{\email}[1]{\def\@email{#1}}
\newcommand{\doi}[1]{\href{https://doi.org/#1}{#1}}

\newcommand{\vol}{00}
\newcommand{\num}{00}
\newcommand{\pubyear}{0000}
\newcommand{\id}{xxxxx}
%%setting page-number
\setcounter{page}{1}


\thispagestyle{empty}
\makeatother

%% Making << \acknowledgment >> command
\newcommand{\acknowledgment}[1]{\vspace{5mm}\singlespacing
    {\noindent\textbf{\textit{Acknowledgment\/}:} #1}
}

%% Setting page numbers
\usepackage{fancyhdr}

%% now reset the headers and footers
\fancyhead{}
\fancyfoot{}
\renewcommand{\headrulewidth}{.3pt}
\voffset = 18pt
\headsep = 3pt

\renewcommand*{\thefootnote}{\fnsymbol{footnote}}

\usepackage[margin=1cm,%
            font=footnotesize,%
            format=hang,%
            labelsep=period,%
            labelfont=bf]{caption}

%%%%%%%%%%%%%%%%%%%%%%%%%%%%%%%%%%%%%%%%%%%%%%%%%%
% PLACE ADDITIONAL PACKAGES HERE (if it is needed):
\usepackage[T1]{fontenc}
\usepackage{lmodern}
\usepackage{mathtools}
\usepackage{array}
\usepackage{enumitem}
\setlist[itemize]{topsep=2pt,itemsep=2pt}
\setlist[enumerate]{topsep=2pt,itemsep=2pt}
\emergencystretch=3em
\hbadness=2000
\setlength{\marginparwidth}{2cm}
%%%%%%%%%%%%%%%%%%%%%%%%%%%%%%%%%%%%%%%%%%%%%%%%%%

%% make the odd pages have the section name on the top right
\fancyhead[RO]{\footnotesize \thepage}
%% make the even pages have the chapter name on the top left
\fancyhead[LE]{\footnotesize \thepage}

%% making page-numbers visible
\pagestyle{fancy}

% Custom macros for the Loyola index family
    \newcommand{\Wuzi}{LO(G;\alpha,\beta,\gamma)}
  \newcommand{\edgepsi}{\psi(x,y;\alpha,\beta,\gamma)}
  \newcommand{\bid}{\mathcal{F}_{\mathrm{BID}}}
% Path to figures
\graphicspath{{figures/}{./}{../figures/}}

% robust figure inclusion: render if present (root or figures/), else a small note
\newcommand{\safefig}[2]{%
  \IfFileExists{#1}{\includegraphics[width=#2\linewidth]{#1}}{%
   \IfFileExists{figures/#1}{\includegraphics[width=#2\linewidth]{#1}}{%
    \fbox{\parbox[c][3cm][c]{0.8\linewidth}{\centering Figure not in project --- upload \texttt{\detokenize{#1}}.}}}}%
}

%%%%%%%%%%%%%%%%%%%%%%%%%%%%%%%%%%%%%%%%%%%%%%%%%%%%%%%%%%%%%%%%%%%%%%%%%%%%%%%%%%%%%%%

\title{Supplementary Material for: The Loyola Index Family: What a Degree-Imbalance Coordinate Can and Cannot Resolve}

\author{Advik Natarajan$^{a}$, Ganesh Shiwakoti$^{b}$, Michael Arockiaraj$^{a}$\footnote{Corresponding author.}}

\address{$^a$Department of Mathematics, Loyola College, Chennai, India\\
    $^b$Department of Mathematics and Statistics, Florida Atlantic University, USA
    }

\email{marockiaraj@gmail.com}

\date{\today}

% !TeX program = pdflatex

\begin{document}

\maketitle
\renewcommand{\thesection}{S\arabic{section}}
\renewcommand{\thetable}{S\arabic{table}}
\renewcommand{\thefigure}{S\arabic{figure}}
\renewcommand{\theequation}{S\arabic{equation}}

\onehalfspacing

\noindent This supplement accompanies the paper ``The Loyola Index
Family: What a Degree-Imbalance Coordinate Can and Cannot Resolve''. Section,
table, figure, equation and proposition numbers without the prefix S
refer to the main paper, and data, scripts and the verifier are available
at \url{https://github.com/Singati2/LOYOLA-PAPER}.

\section{Octane property data: provenance audit and correlation heatmap}\label{sec:s-prov}

Table~4 of the main paper reports the Pearson correlation $r$ between
each index and each of the five physicochemical properties of the
$18$ octane isomers. The properties are the boiling point $T_B$, the
enthalpy of formation $\Delta H_f$, the enthalpy of vaporization
$\Delta H_{\mathrm{vap}}$, the entropy $S$ and the acentric factor
$\omega$. The property data were traced to the NIST
WebBook~\cite{NISTWebBook}, and for each value the provenance file
records the WebBook entry or, where the WebBook lists no value, the
compilation consulted. Acentric factors, which NIST does not list, are
the exception and were taken from KDB-derived compilations.
The same five properties form the octane benchmark used
in~\cite{MovahediGutmanRedzepovicFurtula2026}. $T_B$ is in
$^{\circ}$C, $\Delta H_f$ and $\Delta H_{\mathrm{vap}}$ are in
kcal\,mol$^{-1}$ (the latter at $298$~K), $S$ is in
cal\,mol$^{-1}$K$^{-1}$, and $\omega$ is dimensionless. Values for
$\Delta H_{\mathrm{vap}}$ and entropy $S$ vary slightly across
compilations, and all comparisons here are made within this single,
internally consistent dataset. Calculations use this author-compiled
property table attributed to the NIST WebBook. All $90$
molecule-level values were traced against authoritative sources, with
the NIST WebBook consulted first and KDB-derived compilations used for
the acentric factor, which NIST does not list. The values were
classified under a strict taxonomy. In it, ``verified'' requires
agreement at the value's reported precision, and a converted value
exactly half a display unit away counts as agreement at reported
precision. Of the values, $3$ match exactly, and $29$ match a specific
recorded determination at reported precision after unit conversion,
with the gas phase confirmed for formation enthalpies and entropies.
A further $40$ agree
within the audit tolerances, the source's stated uncertainty or the
span of several determinations without matching any one of them at
reported precision. $10$ lie within explained variation among
authoritative determinations, namely 1947 API Project~44 entries
superseded by later TRC revisions. $7$ could not be verified against a
reachable authoritative value, six of these being gas-phase entropies
absent from the NIST WebBook. For them the 1947 API Project 44 tables
give values differing by $0.07$--$2.2$ units, as they do for isomers
whose later revisions were confirmed. The remaining one is the
$298$\,K vaporization enthalpy of 2,2,3,3-tetramethylbutane. This
compound is a solid at $298$\,K (triple point $374$\,K), and the NIST
WebBook lists standard vaporization enthalpies of $42.94$ and
$42.91$\,kJ\,mol$^{-1}$ (about $10.26$\,kcal\,mol$^{-1}$) together with
sublimation enthalpies. The benchmark value used here,
$8.41$\,kcal\,mol$^{-1}$ ($35.19$\,kJ\,mol$^{-1}$), is that of the
standard octane benchmark compilation and matches no listed
determination. Its physical reference state is therefore unresolved.
Section~S3 reports the $\Delta H_{\mathrm{vap}}$ analysis with this
isomer excluded. One value is a marginal conflict, the acentric factor
of 2,2,3,3-tetramethylbutane ($0.247$ vs.\ compiled $0.251$), while the
acentric factor of 2,2,4-trimethylpentane ($0.305$ vs.\ compiled
$0.303$) agrees only within tolerance. Substituting the alternative
compiled acentric factors changes several of the $\omega$ correlations
in Table~4 of the main paper at the third decimal (e.g.\ $M_2$ moves
from $-0.988$ to $-0.986$). It does not change the property-wise winning
descriptor or any substantive interpretation, and the full sensitivity
computation and the machine-readable provenance trace accompany the
reproducibility package. For 3,3-dimethylhexane,
3-ethyl-3-methylpentane, 2,2,3-trimethylpentane and
2,3,3-trimethylpentane, values of $\Delta H_{\mathrm{vap}}$ that differ
from the NIST WebBook by $0.05$--$0.13$\,kcal\,mol$^{-1}$ circulate in
octane benchmark tables. We use values consistent with NIST
(Appendix~A of the main paper). Two are exact conversions of NIST
determinations ($8.97$, $9.08$), and two lie within the span or stated
uncertainty of the NIST determinations. These are $8.83$, which also
equals the API Project~44 value at the displayed precision, and $8.90$,
which equals the compiled octane dataset value $8.897$ at the displayed
precision. Six gas-phase entropies could not be verified against a
primary source (see above). Four entropies ($n$-octane,
2,2-dimethylhexane, 2,2,4- and 2,3,3-trimethylpentane) differ from
those of a widely used octane QSPR compilation by $0.04$--$2.29$ units,
the $2,3,3$-trimethylpentane value $102.06$ being a compiled-dataset
value (see Section~\ref{sec:s-quality}). Substituting the compilation values
changes all twelve $S$ correlations in Table~4 of the main paper, by up
to $0.043$ in $|r|$, but it does not change the winning descriptor
($HM$). Apart from the four corrected vaporization enthalpies and the
corrected transcription of the 2,3,3-trimethylpentane entropy
(Section~\ref{sec:s-quality}, item~(ii)), the benchmark values are
preserved as supplied for comparability.
\paragraph{Correlation details for Table~4 of the main paper.}
The best correlation per property is split among the reductions and
the pure-$\gamma$ points, with the modified second Zagreb index
$LO(-1, 0, 0)$ strongest on $T_B$ and $\Delta H_f$. The
hyper-Zagreb index $LO(0, 2, 0)$ is strongest on $S$, the second
Zagreb index $LO(1, 0, 0)$ on $\omega$, and the fixed pure-$\gamma$
benchmark point $LO(0, 0, 1)$ on $\Delta H_{\mathrm{vap}}$
($|r| = 0.984$, with $LO(0, 0, 2)$ second at $0.981$). At the
exact-reduction triples the Loyola correlations coincide with the
corresponding classical indices to floating-point precision, as they
must. For each property, a descriptive post-selection paired bootstrap
interval for the observed top-two $|r|$ difference included zero
($B = 10\,000$, resampling the $18$ isomers jointly,
\texttt{default\_rng(1)}). These intervals were not
multiplicity-adjusted across the five properties and do not establish
equality or equivalence. Individual bootstrap half-widths range from
$0.016$ to $0.096$ across the headline cells (percentile bootstrap,
$B = 10\,000$, \texttt{default\_rng(0)}). The value of $LO(0, 0, 1)$ on
$\Delta H_{\mathrm{vap}}$ is that of a fixed benchmark point and
involves no tuning. Recomputing each bolded correlation with single
isomers deleted shifts $|r|$ by at most $0.027$, which indicates that no
single isomer drives the leading correlations.

\paragraph{Descriptor redundancy.} The descriptors are strongly
collinear. This limits what any small $|r|$ difference in
Table~4 of the main paper can mean. Across the twelve compared
descriptors on the $18$ isomers, fourteen descriptor pairs have
$|r| \ge 0.99$, the extreme being $R$ versus $H/2$ at $r = 0.9997$. A
principal component analysis of the standardised $18 \times 12$
descriptor matrix places $90.3\%$ of the variance on the first
component and $99.05\%$ on the first two. The effective rank
(participation ratio of the eigenvalue spectrum) is $1.2$, so the
twelve-index comparison set is essentially a two-dimensional
descriptor space. Differences in the third decimal of $|r|$ between
near-collinear descriptors should not be read as meaningful rankings.
Figure~\ref{fig:octanes} shows the signed-$r$ heatmap of the compared
indices against the five properties, displaying the values of
Table~4 of the main paper.

\begin{figure}[H]
\centering
\safefig{fig_lo_prediction_correlation_heatmap_v3.pdf}{0.9}
\caption{Signed Pearson $r$ between the Loyola reductions, the fixed pure-$\gamma$ benchmark points and the five octane physicochemical properties on the $18$ isomers.}
\label{fig:octanes}
\end{figure}
\section{Effect of parameter tuning near the best indices}\label{sec:tuning}

The benchmark of Section~6.1 of the main paper tunes no parameters
and compares only the exact reductions and the two fixed pure-$\gamma$
benchmark points. It is nonetheless natural to ask whether small,
data-driven adjustments of $(\alpha, \beta, \gamma)$ near the
best-correlating indices could improve the fit, and this section
examines that question as an exploratory check. We take the four anchors
carrying the strongest column correlations in Table~4 of the main paper, namely $M_2 = LO(1, 0, 0)$, $HM = LO(0, 2, 0)$,
${}^{m}\!M_2 = LO(-1, 0, 0)$, and the pure-$\gamma$ point
$LO(0, 0, 1)$. For each anchor we draw $50$ random parameter triples
within $\pm 0.3$ (to two decimal places) of the anchor and correlate
each against the five octane properties. The four pools are drawn, in
the order listed, from one pseudo-random stream (NumPy
\texttt{default\_rng(42)}), each coordinate offset being a uniform draw
on $[-0.3, 0.3]$ rounded to two decimals. The percentile bootstrap uses
\texttt{default\_rng(0)}.
For every anchor--property cell, Table~\ref{tab:lo_tuning} reports the
anchor correlation $|r|_0$, the best of the $50$ perturbations in
sample $|r|_{50}$, the in-sample gain $\Delta_{\mathrm{in}}$ and the
percentile-bootstrap half-width ($B = 10\,000$). It also reports the
leave-one-out correlation of the anchor $\mathrm{LOO}_0$ and a
\emph{nested} leave-one-out correlation $\mathrm{LOO}_{50}$. In the
nested version the best of the $50$ candidates is selected, within each
of the four globally selected anchor pools, using the $17$ training
isomers of each fold only. It is then evaluated on the held-out one.
The resulting out-of-sample gain is $\Delta_{\mathrm{out}}$. The anchors themselves were chosen from the
complete-data correlations of Table~4 of the main paper, and the $18$
pooled fold predictions share heavily overlapping training sets. These
values are therefore descriptive rather than independent replications.

\begin{table}[H]
\centering
\caption{Decimal parameter tuning near the four best-correlating anchors on the $18$ octane isomers ($50$ random triples within $\pm 0.3$ of each anchor). The columns give the anchor correlation $|r|_0$, the best-of-$50$ in sample $|r|_{50}$, the in-sample gain $\Delta_{\mathrm{in}}$ and the percentile-bootstrap half-width $\pm$boot ($B = 10\,000$). They then give the leave-one-out correlation of the anchor $\mathrm{LOO}_0$, the foldwise nested leave-one-out correlation of the tuned family $\mathrm{LOO}_{50}$, computed within the globally selected anchor pool, and the out-of-sample gain $\Delta_{\mathrm{out}} = \mathrm{LOO}_{50} - \mathrm{LOO}_0$. Gains are computed from full-precision values and rounded once, so a displayed gain may differ by one unit in the last place from the difference of the displayed operands.}
\label{tab:lo_tuning}
\scriptsize
\setlength{\tabcolsep}{4pt}
\begin{tabular}{lrrrrrrr}
\toprule
Property & $|r|_0$ & $|r|_{50}$ & $\Delta_{\mathrm{in}}$ & $\pm$boot & $\mathrm{LOO}_0$ & $\mathrm{LOO}_{50}$ & $\Delta_{\mathrm{out}}$\\
\midrule
\multicolumn{8}{l}{\emph{near} $M_2 = LO(1,0,0)$}\\
$T_B$ & $0.497$ & $0.598$ & $+0.100$ & $0.353$ & $0.331$ & $0.458$ & $+0.127$\\
$\Delta H_f$ & $0.542$ & $0.635$ & $+0.093$ & $0.358$ & $0.395$ & $0.540$ & $+0.145$\\
$\Delta H_{\mathrm{vap}}$ & $0.811$ & $0.876$ & $+0.066$ & $0.185$ & $0.768$ & $0.838$ & $+0.071$\\
$S$ & $0.934$ & $0.955$ & $+0.021$ & $0.073$ & $0.909$ & $0.935$ & $+0.025$\\
$\omega$ & $0.988$ & $0.995$ & $+0.007$ & $0.018$ & $0.985$ & $0.985$ & $\phantom{+}0.000$\\
\midrule
\multicolumn{8}{l}{\emph{near} $HM = LO(0,2,0)$}\\
$T_B$ & $0.654$ & $0.740$ & $+0.085$ & $0.257$ & $0.545$ & $0.647$ & $+0.102$\\
$\Delta H_f$ & $0.710$ & $0.786$ & $+0.076$ & $0.255$ & $0.645$ & $0.741$ & $+0.097$\\
$\Delta H_{\mathrm{vap}}$ & $0.903$ & $0.944$ & $+0.041$ & $0.082$ & $0.871$ & $0.921$ & $+0.049$\\
$S$ & $0.971$ & $0.973$ & $+0.002$ & $0.030$ & $0.961$ & $0.957$ & $-0.004$\\
$\omega$ & $0.982$ & $0.989$ & $+0.007$ & $0.025$ & $0.979$ & $0.982$ & $+0.003$\\
\midrule
\multicolumn{8}{l}{\emph{near} ${}^{m}\!M_2 = LO(-1,0,0)$}\\
$T_B$ & $0.855$ & $0.864$ & $+0.009$ & $0.096$ & $0.781$ & $0.759$ & $-0.021$\\
$\Delta H_f$ & $0.891$ & $0.897$ & $+0.006$ & $0.087$ & $0.865$ & $0.837$ & $-0.028$\\
$\Delta H_{\mathrm{vap}}$ & $0.925$ & $0.966$ & $+0.041$ & $0.068$ & $0.900$ & $0.953$ & $+0.053$\\
$S$ & $0.863$ & $0.917$ & $+0.054$ & $0.192$ & $0.810$ & $0.878$ & $+0.068$\\
$\omega$ & $0.794$ & $0.881$ & $+0.088$ & $0.246$ & $0.730$ & $0.853$ & $+0.123$\\
\midrule
\multicolumn{8}{l}{\emph{near} $LO(0,0,1)$}\\
$T_B$ & $0.828$ & $0.835$ & $+0.007$ & $0.161$ & $0.771$ & $0.764$ & $-0.007$\\
$\Delta H_f$ & $0.832$ & $0.842$ & $+0.010$ & $0.152$ & $0.800$ & $0.807$ & $+0.006$\\
$\Delta H_{\mathrm{vap}}$ & $0.984$ & $0.985$ & $+0.001$ & $0.016$ & $0.978$ & $0.978$ & $-0.001$\\
$S$ & $0.928$ & $0.945$ & $+0.017$ & $0.046$ & $0.887$ & $0.915$ & $+0.028$\\
$\omega$ & $0.929$ & $0.945$ & $+0.016$ & $0.074$ & $0.907$ & $0.926$ & $+0.019$\\
\bottomrule
\end{tabular}
\end{table}

Three observations can be made. First, every in-sample gain
$\Delta_{\mathrm{in}}$ is smaller than the corresponding
single-descriptor bootstrap uncertainty scale, although this comparison
is descriptive rather than a formal paired or post-selection test of
the gain. Second, under the foldwise nested evaluation (within the
globally selected anchor pools) these descriptive point-estimate gains
do not vanish. Tuning near $M_2$ raises the out-of-sample correlation
for $T_B$ and $\Delta H_f$ by $+0.13$ and $+0.15$, and tuning near
${}^{m}\!M_2$ raises $\omega$ by $+0.12$. In every such case, however,
the anchor was a \emph{weak} descriptor for that property, and the
tuned correlation only climbs part way toward the level a strong
reduction already attains. Tuning near $M_2$ lifts $T_B$ to $0.46$, whereas
${}^{m}\!M_2$ already reaches $0.78$. Where the anchor is already the
column winner ($M_2$ on $\omega$, $HM$ on $S$, $LO(0, 0, 1)$ on
$\Delta H_{\mathrm{vap}}$), the nested gain for this pool seed is
$-0.0003$, $-0.004$ and $-0.001$. These near-zero values are specific
to the seed. Re-drawing the pools with seeds $0$--$99$, the nested gain
for $M_2$ on $\omega$ is small but positive in most seeds (median
$+0.003$, positive in $97$ of $100$ seeds). The sweep is in the
repository script \texttt{tuning\_pool\_seed\_sweep.py}. Tuning a
column winner therefore changes its nested correlation by at most a few
thousandths. Third, tuning near the pure-$\gamma$ points $LO(0, 0, 1)$
and $LO(0, 0, 2)$ does not exceed the best fixed index under nested
leave-one-out on any of the five properties
(Table~\ref{tab:lo_tuning_bestfixed}). On $\Delta H_{\mathrm{vap}}$
the best fixed index is the pure-$\gamma$ point $LO(0, 0, 1)$ itself.
There the tuned value exceeds it by $0.001$ in sample ($0.985$ vs
$0.984$) and falls $0.001$ below it under leave-one-out (both $0.978$
at three decimals, the differences being those of full-precision
values). These post-selection differences are negligible. On $T_B$,
$\Delta H_f$, $S$, and $\omega$ the best fixed reduction is higher
both in and out of sample.

\begin{table}[H]
\centering
\caption{Decimal tuning near the pure-$\gamma$ points $LO(0,0,1)$ and $LO(0,0,2)$ versus the best fixed index in the tested comparison set for each octane property. The best fixed index is ranked by in-sample $|r|$ (its leave-one-out value in parentheses). The decimal columns give the best over two pools of $50$ random triples within $\pm 0.3$ of $LO(0,0,1)$ and of $LO(0,0,2)$, in sample and under foldwise nested leave-one-out within the globally selected pools. Here the $LO(0,0,1)$ pool is exactly the pool of Table~\ref{tab:lo_tuning}, and the $LO(0,0,2)$ pool is the next $50$ triples of the same \texttt{default\_rng(42)} stream. Indices are ranked by in-sample $|r|$ rather than $\mathrm{LOO}$, as the latter is inflated by degenerate fits for near-zero-correlation indices. Tuning does not exceed the best fixed index under LOO on any property. In sample it exceeds it only on $\Delta H_{\mathrm{vap}}$, by $0.001$, a negligible post-selection difference that does not establish superiority.}
\label{tab:lo_tuning_bestfixed}
\small
\begin{tabular}{lcrrrr}
\toprule
Property & best fixed & $|r|_{\mathrm{fix}}$ & $\mathrm{LOO}_{\mathrm{fix}}$ & $|r|_{\mathrm{dec}}$ & $\mathrm{LOO}_{\mathrm{dec}}$\\
\midrule
$T_B$ & ${}^{m}\!M_2$ & $0.855$ & $0.781$ & $0.835$ & $0.764$\\
$\Delta H_f$ & ${}^{m}\!M_2$ & $0.891$ & $0.865$ & $0.850$ & $0.815$\\
$\Delta H_{\mathrm{vap}}$ & $LO(0,0,1)$ & $0.984$ & $0.978$ & $0.985$ & $0.978$\\
$S$ & $HM$ & $0.971$ & $0.961$ & $0.945$ & $0.915$\\
$\omega$ & $M_2$ & $0.988$ & $0.985$ & $0.945$ & $0.926$\\
\bottomrule
\end{tabular}
\end{table}

In short, local tuning raised several nested leave-one-out point
estimates for weak anchors. Except for the $0.001$ in-sample
$\Delta H_{\mathrm{vap}}$ margin noted above, it did not exceed the
best fixed benchmark in the tested comparison set. No formal
superiority claim is made. All tabled nested values correspond to
positive signed prediction--observation correlations with positive
predictive $Q^{2}$. A constant-descriptor control confirms that
absolute pooled correlation alone can be inflated for degenerate
descriptors (cf.\ the caveat in the caption of
Table~\ref{tab:lo_tuning_bestfixed}).
\section{Matched-budget comparison in the search box $(-2,2)$}\label{sec:s-ablation}

The analyses above tune around fixed anchors and are descriptive.
The question raised by the Gutman--Milovanovi\'c framing is
whether, given identical search budgets, freeing $\gamma$ improves
out-of-sample prediction beyond an optimised parent
$M_{\alpha,\beta}$. This is tested with a fully nested leave-one-out
design. For each property and each outer fold (one isomer held out),
$200$ candidate parameter vectors are drawn once from
$\mathrm{Uniform}(-2,2)$ per coordinate. The candidates are triples
$(\alpha,\beta,\gamma)$ for the Loyola model and pairs $(\alpha,\beta)$
with $\gamma = 0$ for the parent. The budget is identical. Candidates
are rounded to three decimals and drawn from NumPy
\texttt{default\_rng(12345)} for Table~\ref{tab:ablation}, Loyola
triples first. On the $17$ training isomers only, every candidate is
scored by \emph{inner} leave-one-out root-mean-square error of the
one-descriptor linear regression. The best candidate is selected, the
regression is refit on all $17$, and the held-out isomer is predicted.
No quantity computed on the held-out isomer enters selection or fitting
at any point.

\begin{table}[tbp]
\centering
\caption{Matched-budget nested leave-one-out comparison of the
optimised parent $M_{\alpha,\beta}$ (GM) and the optimised Loyola
family ($\gamma$ free) on the five octane properties. Panel (a) is a
\emph{single illustrative candidate stream} (seed $12345$, budget
$200$), where $Q^2$ is the outer predictive coefficient and RMSE the
outer root-mean-square error. ``LO better'' counts outer folds
(of $18$) in which the Loyola model had the strictly smaller absolute
error. Panel (b) summarises the $100$-seed sweep (seeds $0$--$99$),
giving the median $\Delta Q^2$ with its [minimum, maximum], and the
number of seeds in which $\Delta Q^2 > 0$, at budgets $B = 200$ and
$B = 500$. All selection is repeated inside every outer
fold, and all values are full-precision results rounded once.}
\label{tab:ablation}
\small
\setlength{\tabcolsep}{5pt}
\begin{tabular}{lrrrrrr}
\toprule
\multicolumn{7}{l}{\emph{(a) single candidate stream: seed $12345$, budget $200$}}\\
Property & $Q^2_{\mathrm{GM}}$ & $Q^2_{LO}$ & $\Delta Q^2$ & $\mathrm{RMSE}_{\mathrm{GM}}$ & $\mathrm{RMSE}_{LO}$ & LO better\\
\midrule
$T_B$ & $+0.613$ & $+0.748$ & $+0.136$ & $3.700$ & $2.983$ & $10/18$\\
$\Delta H_f$ & $+0.855$ & $+0.549$ & $-0.306$ & $0.463$ & $0.816$ & $4/18$\\
$\Delta H_{\mathrm{vap}}$ & $+0.884$ & $+0.971$ & $+0.086$ & $0.127$ & $0.064$ & $13/18$\\
$S$ & $+0.906$ & $+0.878$ & $-0.029$ & $1.363$ & $1.558$ & $7/18$\\
$\omega$ & $+0.979$ & $+0.956$ & $-0.022$ & $0.005$ & $0.008$ & $5/18$\\
\bottomrule
\end{tabular}

\medskip
\footnotesize
\setlength{\tabcolsep}{4pt}
\begin{tabular}{lrrrr}
\toprule
\multicolumn{5}{l}{\emph{(b) $100$ candidate streams: seeds $0$--$99$}}\\
& \multicolumn{2}{c}{median $\Delta Q^2$ [min, max]} & \multicolumn{2}{c}{seeds with $\Delta Q^2 > 0$}\\
\cmidrule(lr){2-3}\cmidrule(lr){4-5}
Property & $B = 200$ & $B = 500$ & $B = 200$ & $B = 500$\\
\midrule
$T_B$ & $+0.010\,[-0.524,+0.373]$ & $+0.066\,[-0.307,+0.231]$ & $53/100$ & $64/100$\\
$\Delta H_f$ & $-0.120\,[-0.408,+0.206]$ & $-0.095\,[-0.373,+0.118]$ & $14/100$ & $24/100$\\
$\Delta H_{\mathrm{vap}}$ & $+0.069\,[+0.034,+0.102]$ & $+0.068\,[+0.024,+0.102]$ & $100/100$ & $100/100$\\
$S$ & $-0.012\,[-0.140,+0.038]$ & $-0.010\,[-0.098,+0.031]$ & $34/100$ & $31/100$\\
$\omega$ & $-0.010\,[-0.074,+0.016]$ & $-0.006\,[-0.052,+0.012]$ & $26/100$ & $28/100$\\
\bottomrule
\end{tabular}
\end{table}

Once the parent is searched over a box that
does not truncate it, freeing $\gamma$ gives no consistent gain
(Table~5 and Section~6.2 of the main paper). The comparison in the
box $(-2,2)$ is reported first because it shows how the apparent
gain arose, and in that box the result is mixed
(Table~\ref{tab:ablation}). For the illustrative seed, freeing $\gamma$
improves nested predictive performance on two of the five properties
($T_B$, $\Delta Q^2 = +0.136$, and $\Delta H_{\mathrm{vap}}$,
$\Delta Q^2 = +0.086$) and degrades it on three
($\Delta H_f$, $S$, $\omega$). The $\Delta H_f$ degradation is driven
by selection instability. The parent model selected the same
$(\alpha,\beta)$ in all $18$ folds there, while the three-parameter
search selected widely varying triples (fold-to-fold standard
deviation above $1.2$ in both $\alpha$ and $\beta$). Selected
$\gamma$ values are interior to the search box (maximum
$|\gamma| = 0.48$). The parent's selections are not. For
$\Delta H_{\mathrm{vap}}$ the parent chose $(\alpha,\beta) =
(1.021, -1.993)$, on the boundary of the box, in $16$ of the $18$
folds, a choice whose consequences are examined in Section~6.2
of the main paper. On the one property where
$\gamma$ helps most stably, $\Delta H_{\mathrm{vap}}$, the identical
triple $(-0.270, +0.271, +0.239)$ was selected in every fold. With
$n = 18$ molecules these $Q^2$ differences carry substantial
selection noise. We attach no inferential claim to them.

Because the comparison depends on a random candidate stream, the full
nested design was repeated for one hundred seeds ($0$--$99$) at both
budgets ($200$ and $500$ candidates per model). This makes $1000$
configurations in all. A paired control was run alongside, comparing
the Loyola candidates against the \emph{same} $(\alpha,\beta)$
coordinates with $\gamma$ set to zero, which isolates the effect of
$\gamma$ from the independent candidate streams. Freeing
$\gamma$ improved nested predictive performance for
$\Delta H_{\mathrm{vap}}$ in every seed of this sweep ($100/100$ at
both budgets, and $100/100$ in the paired control). The median was
$\Delta Q^2 = +0.069$ and $+0.068$, with range $+0.024$ to $+0.102$.
It \emph{usually} degraded performance for $\Delta H_f$ ($86/100$ and
$76/100$ seeds), while for $T_B$ the split is close to even at budget
$200$ ($53/100$) but leans towards improvement at budget $500$
($64/100$). $S$ and $\omega$ usually degrade slightly. These counts
describe optimiser variance over random candidate streams on one fixed
set of $18$ molecules, which are never resampled. The counts therefore
carry no inferential weight. In particular, $100/100$ does not establish that
$\gamma$ improves $\Delta H_{\mathrm{vap}}$ prediction for other
molecules or data sets. The result is also sensitive to a single data
cell, the $\Delta H_{\mathrm{vap}}$ value of 2,2,3,3-tetramethylbutane,
whose physical reference state is unresolved (Section~S1). The
complete nested comparison was repeated with that isomer excluded
($17$ molecules, same seeds, budgets and candidate streams). Freeing
$\gamma$ then improved $\Delta H_{\mathrm{vap}}$ in $93/100$ seeds at
budget $200$ and $96/100$ at budget $500$ (paired control $95/100$ and
$96/100$). The median was $\Delta Q^2 = +0.046$ and $+0.051$, with
range $-0.048$ to $+0.098$. The gain survives the exclusion but is no
longer sign-invariant.
These comparisons share the box $(-2,2)$, which truncates the parent
along the direction $\beta \approx -2\alpha$. Section~6.2
and Table~5 of the main paper show that the
vaporization-enthalpy gain shrinks to about zero once the parent is
searched over a box that does not truncate it.

\section{Nonane replication test: protocol, exploratory transfer and
molecule bootstrap}\label{sec:s-external}

Two limitations of the octane analysis are its size ($18$ molecules) and
its restriction to index models. Both were addressed under a protocol
recorded in a git commit of the reproducibility package before any
nonane value was collected. The identical nested design was applied to
the $35$ constitutional nonane isomers. Boiling points ($35$ values) and
standard vaporization enthalpies ($34$ values, none being listed for
3-ethyl-4-methylhexane) were taken from the NIST
WebBook~\cite{NISTWebBook}, with every value traced to its archived
source page. The two properties differ in evidential status. Twenty of
the $35$ boiling points are NIST averages of experimental
determinations ($13$ quoted only to $1$\,K). Five are single NIST
determinations and ten are the most recent primary determination listed
by NIST (Deviation~D1 of the protocol). All are experimental
measurements, so the $T_B$ test is experimentally grounded.
Of the vaporization enthalpies, only $4$ of $34$ are calorimetric
determinations and one is a NIST average. The other $29$ come from two
compilations. These are the 1961 tabulation of
Labauf, Greenshields and Rossini~\cite{LabaufGreenshieldsRossini1961},
which covers all $35$ nonanes and $75$ decanes and is not a set of
individual calorimetric measurements, and a 1971 vapour-pressure
handbook~\cite{WilhoitZwolinski1971}. Such tabulations may rest
partly on structural-increment estimates, so agreement with them
measures consistency with the compiled values rather than independent
experimental validation. A degree-pair-based predictor may also be
favoured if the compiled values were themselves derived from structural
increments. The pre-registered
decision rule called the $\gamma$ effect \emph{replicated} if the Loyola
model beat the parent in at least $80$ of $100$ seeds at both budgets
with positive median $\Delta Q^2$. Four baselines were added, each
fitted entirely within the outer training fold. They are the
training-fold mean, the best of the ten classical indices of Table~4 of the main paper, and the irregularity index $IRLA$. The fourth is ridge
regression on the full edge-degree-pair count vector, with its penalty
chosen by inner leave-one-out, and the resulting $Q^2$ values are in
Table~6 of the main paper.

On the nonane replication set the pre-registered verdicts are as
follows. For $\Delta H_{\mathrm{vap}}$ (compilation values) the $\gamma$
effect met the pre-registered optimiser-stability criterion and
\emph{replicates}, LO being better than GM in $95/100$ and $98/100$
seeds, with median $\Delta Q^2 = +0.048$ and $+0.053$. For $T_B$ the
verdict is \emph{mixed} ($58/100$ and $63/100$). It is unchanged in
direction under the pre-registered sensitivity analysis D4 of the
protocol. That analysis replaces the NIST average for
2,3,3,4-tetramethylpentane, $413 \pm 6$~K, by the median of all eight
listed determinations including a 1989 handbook value and a 1939
outlier, $414.62$~K. LO is then better in $60/100$ and $66/100$ seeds,
with median $\Delta Q^2 = +0.041$ and $+0.052$. The pre-registered
protocol fixed the box of half-width $2$, which, as
Table~5 of the main paper shows, binds the parent. In a box that does not, the nonane gain
disappears (median $\Delta Q^2 = -0.002$ under the hybrid control), and
the pre-registered verdict is reported as obtained, with this deviation
recorded. The pattern behind the vaporization result is the same on
both sets. Models whose edge weight depends on imbalance only through
$q^2$ inside the box of half-width $2$ (the classical indices,
including $GA = \sum_{uv}\sqrt{1-q_{uv}^2}$, and GM) reach
$Q^2 = 0.853$--$0.888$. Models with a term linear in $q$ or with
unrestricted counts ($IRLA$, ridge, $LO(0,0,1)$, LO) reach
$Q^2 = 0.897$--$0.957$. This split holds only inside the narrow box. In
the box of half-width $12$ the parent alone reaches $0.942$ and $0.892$
(main paper, Table~5). The fixed point $LO(0,0,1)$ involves no tuning
at all, and it matches or beats the tuned Loyola model on both
sets ($0.957$ and $0.903$ against medians $0.953$ and $0.900$). It is
also the index selected in every fold for $\Delta H_{\mathrm{vap}}$
when admitted to the fixed-index comparison. It is closely related to
$IRLA$. One has
$LO(G;0,0,\gamma) = m + \tfrac{\gamma}{2}\,IRLA(G) + O(\gamma^2)$ as
$\gamma \to 0$, and on both sets the correlation between $LO(0,0,1)$
and $IRLA$ is $0.996$. The baselines, however, also show the limit of
the claim. Ridge regression on the full degree-pair count vector
predicts the nonanes better than every index model, for both
properties ($Q^2 = 0.927$ for $\Delta H_{\mathrm{vap}}$, $0.682$ for
$T_B$). On the octanes it matches the Loyola model on
$\Delta H_{\mathrm{vap}}$ and is weaker elsewhere. This is consistent
with Proposition~12 of the main paper and the remark after
Proposition~7. The imbalance coordinate re-weights degree-pair
information rather than adding to it, and one interpretable descriptor
recovers most of what a model on all ten degree-pair counts achieves
for the vaporization enthalpy. It does not outperform such a model. The
specific exponential form is not essential either, since the
parameter-free aggregate $IRLA$ comes within $0.003$ of the median
Loyola model on the nonanes ($0.897$ against $0.900$). For this
property a term linear in the imbalance accounts for the effect.

Finally, the nested design refits every model on the nonanes. It
therefore tests whether the \emph{procedure} replicates rather than
whether an octane-trained model transfers. An exploratory analysis
followed. It was not pre-registered, and it was run in the box of
half-width $2$, so that it inherits the truncation of the parent. In
it, the parameters selected on all $18$ octanes were frozen and applied
to the nonanes. With the parameters frozen and only the linear
calibration refitted, the Loyola model beat the parent for
$\Delta H_{\mathrm{vap}}$ in $100/100$ seeds at both budgets (median
$Q^2 = 0.910$ and $0.916$ against $0.864$ and $0.865$). It beat the
parent for $T_B$ in $65/100$ and $82/100$. With the octane calibration
frozen as well, both models extrapolate poorly from $C_8$ to $C_9$. The
Loyola model is then worse than the parent (median $Q^2$ between
$-2.22$ and $-0.72$ for LO against $0.06$ to $0.44$ for GM). Because
both comparisons inherit the truncated box, they do not bear on whether
$\gamma$ adds predictive information.

The seed counts above measure optimiser variance only. To indicate
molecule-level uncertainty we also computed, post hoc, a paired
bootstrap over molecules ($2000$ resamples) of the median-over-seeds
$\Delta Q^2$, conditional on the fitted outer predictions. Selection is
not repeated within resamples. The intervals therefore do not account
for uncertainty from refitting and model selection. The $95\%$ intervals
for LO minus GM on $\Delta H_{\mathrm{vap}}$ are $[+0.014, +0.111]$ on
the octanes but $[-0.006, +0.113]$ on the nonanes at budget $200$
($[-0.005, +0.119]$ at $500$). Although the nonane gain met the
pre-registered criterion, its interval narrowly includes zero, as does
the interval for LO minus ridge on the nonanes, $[-0.101, +0.032]$.
For $T_B$ the intervals are wide ($[-0.212, +0.222]$ for LO minus GM on
the nonanes). The one octane interval that excludes zero in the other
direction is that for $\Delta H_f$, where freeing $\gamma$ is worse
($[-0.281, -0.044]$). All these intervals refer to the box of
half-width $2$. Given the search-box sensitivity of Table~5 of the main paper, no $\gamma$ effect
remains to be established. The intervals are reported for
completeness.
\section{Degeneracy across tree orders}\label{sec:s-degeneracy}

Degeneracy~\cite{Konstantinova1996Degeneracy} measures the failure of
an index to discriminate non-isomorphic graphs, so lower values are
better. On the $106$ non-isomorphic
trees of order $10$ we use the distinct-value convention
\[
\mathrm{Degeneracy}(\%) = 100\Bigl(1 - \tfrac{k}{N}\Bigr),
\quad k = \#\{\text{distinct index values}\},\ N = 106,
\]
i.e.\ the number $N - k$ of structures in excess of the distinct
values, as a percentage of $N$. This is not the share of structures
lacking a unique value, which is larger. For $LO(0,0,1)$ on this set,
$45$ of the $106$ trees ($42.5\%$) share their value with another
tree. The pure-$\gamma$ points $LO(0, 0, 1)$ and
$LO(0, 0, 2)$ attain degeneracy $25.5\%$ ($79$ distinct values),
among the lowest of all compared indices and tied with $R$, $GA$
and $AG$. The $106$ trees realise only $79$ distinct
edge-degree-pair count vectors, so $25.5\%$ is the minimum degeneracy
attainable by \emph{any} degree-pair-based index on this set, and the
pure-$\gamma$ points attain this floor. Theorem~2 of the main paper
gives the reason and the order at which this stops. The first
Zagreb index $LO(0, 1, 0)$ ($83.0\%$) and the second Zagreb index
$LO(1, 0, 0)$ ($67.0\%$) are the most degenerate
(Figure~\ref{fig:degeneracy}).

\begin{figure}[H]
\centering
\safefig{fig_lo_degeneracy_order10_trees_v2.pdf}{0.9}
\caption{Degeneracy percentage of the compared indices on the $106$ non-isomorphic trees of order $10$. The dashed line marks the degree-pair-profile floor, $25.5\% = 100(1 - 79/106)$. The $106$ trees realise only $79$ distinct edge-degree-pair count vectors, so no vertex-degree-based index can fall below it.}
\label{fig:degeneracy}
\end{figure}

To test whether this floor attainment is peculiar to order $10$, the
same computation was repeated for all non-isomorphic trees of orders
$7$ through $12$. It was also repeated, separately, for the molecular
subsets with $\Delta \le 4$ (Table~\ref{tab:multiorder}), and the
pure-$\gamma$ points attain the degree-pair-profile floor at
\emph{every} order in both sets. The Randi\'c index also attains the
floor at every order in this range, while $\chi$ falls off the floor
from order $10$. Floor attainment therefore distinguishes the
pure-$\gamma$ points from the most degenerate classical indices
($M_1$, $M_2$) but not from $R$. The imbalance weights therefore lose nothing, although they are not
unique in this respect.
Theorem~2 of the main paper turns this into an exact statement. The
pure-$\gamma$ points separate precisely the $q$-histograms, which
coincide with the degree-pair profiles up to order $12$ (order $15$ for
$\Delta \le 4$) and not beyond. Among the trees of order $13$ the
profiles
$\{(1,2)^{1}, (1,3)^{2}, (1,6)^{5}, (2,3)^{2}, (2,6)^{1}, (3,3)^{1}\}$
and $\{(1,3)^{3}, (1,6)^{5}, (2,2)^{1}, (2,3)^{2}, (3,6)^{1}\}$
(exponents are edge multiplicities) share the $q$-histogram
$\{0^{1}, \tfrac15^{2}, \tfrac13^{1}, \tfrac12^{3}, \tfrac57^{5}\}$,
and there the pure-$\gamma$ points realise $566$ distinct values against
$570$ profiles. Among molecular trees ($\Delta \le 4$) the first
failure occurs at order $16$ ($809$ values against $810$ profiles).
The Randi\'c index still attains the floor at order $13$ ($570$ values)
and falls below it at order $14$ ($1095$ values against $1100$
profiles, where the pure-$\gamma$ points give $1090$).

\begin{table}[H]
\centering
\caption{Discrimination across tree orders $n = 7$--$12$. The columns
give the number of trees $N$, the number of distinct edge-degree-pair
profiles (the floor attainable by any vertex-degree-based index), and
distinct-value counts for representative indices. ``mol.''\ denotes the
molecular subset $\Delta \le 4$.}
\label{tab:multiorder}
\small
\begin{tabular}{llrrrrrrr}
\toprule
$n$ & set & $N$ & floor & $M_1$ & $M_2$ & $R$ & $LO(0,0,1)$ & $LO(0,0,2)$\\
\midrule
$7$ & all & $11$ & $11$ & $7$ & $9$ & $11$ & $11$ & $11$\\
$7$ & mol. & $9$ & $9$ & $5$ & $7$ & $9$ & $9$ & $9$\\
$8$ & all & $23$ & $21$ & $9$ & $17$ & $21$ & $21$ & $21$\\
$8$ & mol. & $18$ & $16$ & $6$ & $13$ & $16$ & $16$ & $16$\\
$9$ & all & $47$ & $40$ & $13$ & $22$ & $40$ & $40$ & $40$\\
$9$ & mol. & $35$ & $28$ & $7$ & $16$ & $28$ & $28$ & $28$\\
$10$ & all & $106$ & $79$ & $18$ & $35$ & $79$ & $79$ & $79$\\
$10$ & mol. & $75$ & $49$ & $8$ & $21$ & $49$ & $49$ & $49$\\
$11$ & all & $235$ & $152$ & $21$ & $43$ & $152$ & $152$ & $152$\\
$11$ & mol. & $159$ & $83$ & $9$ & $25$ & $83$ & $83$ & $83$\\
$12$ & all & $551$ & $294$ & $27$ & $57$ & $294$ & $294$ & $294$\\
$12$ & mol. & $355$ & $137$ & $10$ & $28$ & $137$ & $137$ & $137$\\
\bottomrule
\end{tabular}
\end{table}

\paragraph{Fixed fractional exponents (exploratory).} A script in the
repository (directory \texttt{structural}) evaluates the nine fixed
points
\begin{multline*}
(\alpha,\beta) \in \bigl\{(0,\tfrac12), (\tfrac12,0), (0,-\tfrac12),
(-\tfrac12,0), (\tfrac12,-1),\\
(-\tfrac12,1), (\tfrac13,0), (0,\tfrac13), (\tfrac16,0)\bigr\}
\end{multline*}
with $\gamma = 0$ and $\gamma = 1$. We first consider discrimination on
all trees of orders $12$--$16$, where every point off the ratio plane
with $\gamma = 1$ reaches or nearly reaches the profile floor. At order
$16$, for example, the counts are $4069$, $4067$, $4069$, $4068$,
$4069$ and $4069$ of $4069$ for $(0,\tfrac12)$, $(\tfrac12,0)$,
$(0,-\tfrac12)$, $(-\tfrac12,0)$, $(\tfrac13,0)$, $(0,\tfrac13)$. With
$\gamma = 0$ the same points give $3475$, $3940$, $3476$, $4007$,
$4009$ and $3476$, while the two points on the ratio plane give $3996$
for both values of $\gamma$. We next consider prediction (fixed points,
leave-one-out $Q^2$, no tuning). Here $\gamma = 1$ beats $\gamma = 0$
in $52$ of $81$ point--dataset combinations, mostly for the
vaporization enthalpy and decane boiling points, and loses mostly on
the octane entropy and acentric factor. Every fixed fractional point
stays at $Q^2 \le 0.21$ on the decanes. The full output is committed
with the script.

\section{Structure sensitivity: protocol and published control}\label{sec:s-ss}

The structure-sensitivity criterion for descriptor quality was
introduced by Furtula, Gutman and
Dehmer~\cite{FurtulaGutmanDehmer2013SS}, and a fingerprint-based
variant is due to Raki\'c and Furtula~\cite{RakicFurtula2019SS}. We
implement the published protocol on the $75$ decane isomer trees
$\Psi$. For each tree $G$, the similar-structure set
$S(G) \subset \Psi$ consists of the non-isomorphic trees obtainable
from $G$ by relocating a single edge within the molecular class
(graph edit distance two, $\Delta \le 4$ preserved). Then
\begin{gather*}
SS(TI, G) = \frac{1}{|S(G)|} \sum_{\Omega \in S(G)}
\frac{|TI(\Omega) - TI(G)|}{TI(G)},\\
Abr(TI, G) = \max_{\Omega \in S(G)}
\frac{|TI(\Omega) - TI(G)|}{TI(G)},
\end{gather*}
and the class-level $SS(TI)$ and $Abr(TI)$ are the means over all
$75$ trees. On this class every $S(G)$ is nonempty (sizes $7$ to
$43$). The neighbour relation is symmetric, and $Abr \ge SS$ holds
throughout, as it must. As an external control, the implementation
reproduces the published decane-class values of Barman and
Das~\cite{BarmanDas2026HSO} for the seven indices common to both
studies ($M_1$, ${}^{m}\!M_2$, $R$, $\chi$, $H$, $GA$, $AG$). These
give fourteen tabulated $SS$/$Abr$ values, of which eleven agree exactly
at the four displayed decimals and the remaining three differ by one
unit in the fourth decimal. The cell-by-cell comparison is in the file
\texttt{fgd\_published\_control.csv} of the reproducibility package,
while our values for all twelve compared indices are in
Table~3 of the main paper. The published desideratum is $SS$ large
and $Abr$ small~\cite{FurtulaGutmanDehmer2013SS}. Under it, the
pure-$\gamma$ point $LO(0,0,2)$ has a structure sensitivity
indistinguishable at three decimals from that of the hyper-Zagreb
index ($0.1257$ against $0.1256$) at a
markedly smaller abruptness ($0.274$ against $0.313$). The
sensitivity-to-abruptness ratio is our own summary of the two published
quantities, not part of the published criterion. The ratios of the
pure-$\gamma$ points
($0.4579$ and $0.4565$) are comparable to those of $H/2$, $GA$ and
$R$ ($0.448$--$0.451$). In view of the redundancy analysis of
Section~6.1 of the main paper and Section~S1, these third-decimal
differences are not a ranking. Two qualifications limit this. First,
$SS$ is not a fixed property of the family, since along the pure-$\gamma$
line it grows with $\gamma$ ($0.031$, $0.062$, $0.126$, $0.319$ at
$\gamma = 0.5, 1, 2, 5$). The ratio peaks near $\gamma = 2$ and then
falls ($0.427$ at $\gamma = 5$). Second, the parameter-free aggregate
$IRLA$, which the pure-$\gamma$ points approximate to first order in
$\gamma$, has a higher ratio ($SS = 0.2343$, $Abr = 0.5066$, ratio
$0.4625$). The favourable sensitivity--abruptness balance is
therefore a property of linear imbalance weighting rather than of the
exponential form.

\section{Lean 4 formalisation: scope and encoding}\label{sec:s-lean}

The directory \texttt{formal/} of the repository contains a Lean~4
formalisation (Lean v4.34.1, Mathlib v4.34.1) of the following
statements of the main paper. \texttt{formal/verify.sh} builds it and
checks that no \texttt{sorry}, \texttt{axiom} or \texttt{native\_decide}
occurs, and that every theorem depends only on the standard
axioms (\texttt{propext}, \texttt{Classical.choice}, \texttt{Quot.sound}).

\emph{Encoding.} Graphs are not formalised as Mathlib graphs, and an edge
set is a finite family of edges with a map to (positive real) degree
pairs, or a multiset of degree pairs. Hypotheses such as ``maximum degree
$\Delta$'' are stated as bounds on these pairs, so the link ``these are
the vertex degrees of a graph'' is not formalised. The only
exception is the two order-$13$ trees below, which are realised by
explicit parent arrays.

\emph{Formalized} (inequalities and identities, with their hypotheses).
The formalised statements are the upper and lower $(\Delta,\delta)$
bounds, the Cauchy--Schwarz interpolation bound and the ratio bound
(and the bracket form). They also include the convexity of
$\theta \mapsto \log LO(G;\theta)$ for any nonempty edge set (via
H\"older), and the spanning statement of the independence lemma for
every $\Delta \ge 3$ (via $\det W = \tfrac16\ln 2\,(8\ln 3 - 13\ln 2)
\neq 0$, using $3^8 \neq 2^{13}$). Finally, the identity
$1 - q^2 = 4xy/(x+y)^2$ is formalised, as is the order-$13$ profile
pair with equal $q$-histograms and hence equal pure-$\gamma$ values
for every real $\gamma$, together with explicit $13$-vertex trees
realising both profiles.

\emph{Not formalised.} The equality characterisations are not
formalised, and neither are the gradient and Hessian formulas and strict
convexity. The ratio-plane theorem, the $\gamma$-crossing bound and
the identifiability proposition are not formalised either. These rely
on the Lindemann--Weierstrass theorem, Laguerre's rule and the theory
of exponential families, which are not available in Mathlib. All
statistical and chemometric results are outside the formalisation, as
are the closed forms and the
ceiling and generic-discrimination propositions. The identifiability
consequence of Lemma~1 together with its $\Delta = 1, 2$ cases is
also not formalised. A theorem-by-theorem map is in
\texttt{formal/README.md}.

\section{Further directions}\label{sec:s-further}

The following open problems remain:
\begin{enumerate}
\item Sharp two-sided bounds for $\Wuzi$ via degree-based indices
lying outside the family, such as the atom-bond-connectivity index $ABC$ and the
Sombor index.


\item Extremal-graph characterisations (not attempted in this paper) for
several graph classes, including trees, unicyclic and bicyclic graphs.

\item Chemometric extension to additional chemical compound classes such as benzene derivatives, alkane homologues beyond the octanes, nonanes and
decanes studied here and to compound-class-specific QSPR endpoints.
\end{enumerate}

The first orders at which the pure-$\gamma$ points, the Randi\'c index
and $GA$ lose discrimination are summarised in Table~2 of the main paper.



\section{Data quality and open provenance items}\label{sec:s-quality}

Every value in Table~\ref{tab:octane-data} was traced to a source
(per-value record in the repository). The items that remain
open are listed here so that they can be read in one place.
\begin{enumerate}[label=(\roman*)]
\item \emph{Thirteen gas-phase entropies have no $S^{\circ}$(gas) entry
  in the NIST WebBook.} Seven of them (2,3-, 2,4-, 2,5- and
  3,3-dimethylhexane, 3-ethylhexane, 2,2,3- and 2,2,4-trimethylpentane)
  agree, within tolerance or as documented revisions between editions,
  with the 1947 API Research Project~44 tables or a compiled octane
  dataset. The remaining six (4-methylheptane, 2,2-dimethylhexane,
  3,4-dimethylhexane, 3-ethyl-2-methylpentane, 3-ethyl-3-methylpentane,
  2,3,3-trimethylpentane) could not be matched to any primary source. The 1947 API Research Project 44
  tables~\cite{APIProject44Circ461} list gas-phase entropies for all six.
  These differ from the values used by $0.07$--$2.2$\,cal\,mol$^{-1}$\,K$^{-1}$,
  as that edition does for isomers whose later revisions we could
  confirm. No later primary determination was located, so the values used
  are compiled-dataset values. Substituting the 1947 values changes the
  $S$ correlations by at most $0.018$ in $|r|$ and puts $M_1$ ahead of
  $HM$ by $0.004$.
\item \emph{2,3,3-trimethylpentane $S$.} The value $101.31$ found in our
  earlier compilation is a copy of the 2,2,3-trimethylpentane entry and is
  not used. The value used, $102.06$, is that of the compiled
  $18$-octane dataset reproduced by Mondal et al.~\cite{MondalDePal2019QSPR}
  (the one-decimal compilation of Ediz~\cite{Ediz2017Octane} gives
  $102.10$). This change moves the $S$ correlations by at most $0.009$, and the
  entry used remains a secondary value (item~(i)).
\item \emph{Four entropies differ from the compiled octane dataset of
  Ediz~\cite{Ediz2017Octane}} (octane $111.55$ vs $111.70$;
  2,2-dimethylhexane $103.13$ vs $103.40$; 2,2,4-trimethylpentane $101.81$
  vs $104.10$; 2,3,3-trimethylpentane $102.06$ vs $102.10$). For octane
  the value used agrees with the NIST WebBook ($111.63$). Substituting the
  compilation values changes all twelve $S$ correlations, by up to $0.043$
  in $|r|$, and moves the margin between $HM$ and $M_1$ at the top from
  $0.001$ to $0.007$.
\item \emph{2,2,3,3-tetramethylbutane $\Delta H_{\mathrm{vap}}$} is a
  benchmark-compilation value ($8.41$) whose reference state is unresolved.
  The compound is a solid at $298$~K, and NIST lists sublimation and standard
  vaporization enthalpies near $10.3$~kcal\,mol$^{-1}$. We repeated the
  matched-budget comparison of Supplementary Section~S3 without this
  isomer, and the vaporization-enthalpy gain in the $(-2,2)$ box persists
  ($93/100$ and $96/100$ seeds instead of $100/100$). The
  $\Delta H_{\mathrm{vap}}$ column winner of Table~4 of the main paper is
  unchanged, and no conclusion of the paper rests on this value.
\item \emph{Acentric factors} of 2,2,4-trimethyl\-pentane and
  2,2,3,3-tetra\-methyl\-butane differ from an alternative compilation by
  $0.002$--$0.004$, and these alternative values change several $\omega$
  correlations at the third decimal and leave the $\omega$ ranking
  unchanged.
\end{enumerate}
None of these items moves a correlation of Table~4 of the main paper by
more than $0.043$. The only bolded entry at stake is the $S$ winner,
where $HM$ leads $M_1$ by $0.001$ with the values used. The 1947
substitution of item~(i) reverses that order by $0.004$, a margin
without meaning at this resolution. None of the items enters the
pre-registered nonane and boiling-point tests of Sections~6.3 and~6.4 of the main paper.
Those tests use no entropy, vaporization-enthalpy or acentric-factor value of
Table~\ref{tab:octane-data}, and the pooled boiling-point test re-extracts
the octane boiling points under its own rule. Item~(ii) was a transcription error in our
earlier data file, corrected here, while the remaining items are limitations
of the compiled octane benchmark rather than of the analysis.



Table~\ref{tab:octane-data} lists the full $18$-isomer dataset used
in Section~6 of the main paper, with $T_B$ in $^{\circ}$C, $\Delta H_f$ and
$\Delta H_{\mathrm{vap}}$ in kcal\,mol$^{-1}$
($\Delta H_{\mathrm{vap}}$ at $298$~K), $S$ in
cal\,mol$^{-1}$K$^{-1}$ and $\omega$ dimensionless. The
$\Delta H_{\mathrm{vap}}$ values of 3,3-dimethylhexane ($8.97$),
3-ethyl-3-methylpentane ($9.08$), 2,2,3-trimethylpentane ($8.83$) and
2,3,3-trimethylpentane ($8.90$) are values consistent with the NIST
WebBook determinations, and Section~\ref{sec:s-prov} gives the
exact-versus-within-uncertainty status of each. The
$\Delta H_{\mathrm{vap}}$ entry of 2,2,3,3-tetramethylbutane is the
benchmark-compilation value discussed above. The $S$ value of
2,3,3-trimethylpentane ($102.06$) is a compiled-dataset value
(this section).

Table~\ref{tab:octane-data} has several open provenance items, namely
the unverified entropies, the vaporization enthalpy of
2,2,3,3-tetramethylbutane, the corrected transcription of the
2,3,3-trimethylpentane entropy and the alternative compiled values.
These items, with their effect on Table~\ref{tab:octane-data}, are
collected in this section, and none of them enters the
pre-registered tests of Sections~6.3 and~6.4 of the main paper.

\begin{table}[H]
\centering\small
\caption{Physicochemical properties of the $18$ octane isomers.}
\label{tab:octane-data}
\begin{tabular}{lrrrrr}
\toprule
Isomer & $T_B$ & $\Delta H_f$ & $\Delta H_{\mathrm{vap}}$ & $S$ & $\omega$\\
\midrule
$n$-octane & 125.6 & $-49.82$ & 9.92 & 111.55 & 0.398\\
2-methylheptane & 117.6 & $-51.50$ & 9.48 & 109.84 & 0.378\\
3-methylheptane & 118.9 & $-50.82$ & 9.52 & 111.26 & 0.371\\
4-methylheptane & 117.7 & $-50.69$ & 9.48 & 109.32 & 0.372\\
2,2-dimethylhexane & 106.8 & $-53.71$ & 8.92 & 103.13 & 0.339\\
2,3-dimethylhexane & 115.6 & $-51.13$ & 9.27 & 108.02 & 0.348\\
2,4-dimethylhexane & 109.4 & $-52.44$ & 9.03 & 106.98 & 0.344\\
2,5-dimethylhexane & 109.1 & $-53.21$ & 9.05 & 105.72 & 0.357\\
3,3-dimethylhexane & 111.9 & $-52.61$ & 8.97 & 104.74 & 0.322\\
3,4-dimethylhexane & 117.7 & $-50.91$ & 9.32 & 106.59 & 0.340\\
3-ethyl-2-methylpentane & 115.6 & $-50.48$ & 9.21 & 106.06 & 0.330\\
3-ethyl-3-methylpentane & 118.3 & $-51.38$ & 9.08 & 101.48 & 0.302\\
3-ethylhexane & 118.5 & $-50.40$ & 9.48 & 109.43 & 0.362\\
2,2,3-trimethylpentane & 109.8 & $-52.61$ & 8.83 & 101.31 & 0.300\\
2,2,4-trimethylpentane & 99.2 & $-53.57$ & 8.40 & 101.81 & 0.305\\
2,3,3-trimethylpentane & 114.8 & $-51.73$ & 8.90 & 102.06 & 0.291\\
2,3,4-trimethylpentane & 113.5 & $-51.97$ & 9.01 & 102.39 & 0.317\\
2,2,3,3-tetramethylbutane & 106.5 & $-53.99$ & 8.41 & 93.06 & 0.247\\
\bottomrule
\end{tabular}
\end{table}

\singlespacing
\begin{thebibliography}{00}

  \bibitem{BarmanDas2026HSO} J. Barman, S. Das, Geometric approach to
  degree-based topological index: Hyperbolic Sombor index, \textit{MATCH
  Commun. Math. Comput. Chem.\/} \textbf{95} (2026) 63--94.

  \bibitem{FurtulaGutmanDehmer2013SS} B. Furtula, I. Gutman, M. Dehmer, On
  structure-sensitivity of degree-based topological indices, \textit{Appl.
  Math. Comput.\/} \textbf{219} (2013) 8973--8978.

  \bibitem{Konstantinova1996Degeneracy} E. V. Konstantinova, The
  discrimination ability of some topological and information distance indices
  for graphs of unbranched hexagonal systems, \textit{J. Chem. Inf. Comput.
  Sci.\/} \textbf{36} (1996) 54--57.

  \bibitem{LabaufGreenshieldsRossini1961} A. Labbauf, J. B. Greenshields,
  F. D. Rossini, Heats of formation, combustion, and vaporization of the
  35 nonanes and 75 decanes, \textit{J. Chem. Eng. Data\/} \textbf{6}
  (1961) 261--263.

  \bibitem{NISTWebBook} P. J. Linstrom, W. G. Mallard (Eds.), \textit{NIST
  Chemistry WebBook, NIST Standard Reference Database Number 69\/}, National
  Institute of Standards and Technology, Gaithersburg, 2025,
  https://doi.org/10.18434/T4D303.

  \bibitem{MovahediGutmanRedzepovicFurtula2026} F. Movahedi, I. Gutman, I.
  Red\v{z}epovi\'c, B. Furtula, Diminished Sombor index, \textit{MATCH Commun.
  Math. Comput. Chem.\/} \textbf{95} (2026) 141--162.

  \bibitem{RakicFurtula2019SS} M. Raki\'c, B. Furtula, A novel method for
  measuring the structure sensitivity of molecular descriptors, \textit{J.
  Chemometrics\/} \textbf{33} (2019) \#e3138.

  \bibitem{WilhoitZwolinski1971} R. C. Wilhoit, B. J. Zwolinski,
  \textit{Handbook of Vapor Pressures and Heats of Vaporization of
  Hydrocarbons and Related Compounds\/}, Texas A\&M Research Foundation,
  College Station, 1971.

  \bibitem{APIProject44Circ461} F. D. Rossini, K. S. Pitzer, W. J. Taylor,
  J. P. Ebert, J. E. Kilpatrick, C. W. Beckett, M. G. Williams, H. G. Werner,
  \textit{Selected Values of Properties of Hydrocarbons\/}, Natl. Bur.
  Stand. Circ. 461, U.S. Government Printing Office, Washington, 1947.

  \bibitem{Ediz2017Octane}
  S. Ediz, Predicting some physicochemical properties of octane isomers:
  a topological approach using ev-degree and ve-degree Zagreb indices,
  arXiv:1701.02859 (2017), \url{https://arxiv.org/abs/1701.02859}.

  \bibitem{MondalDePal2019QSPR}
  S. Mondal, A. Dey, N. De, A. Pal, QSPR analysis of some novel
  neighbourhood degree-based topological descriptors, \textit{Complex
  Intell. Syst.\/} \textbf{7} (2021) 977--996,
  \url{https://doi.org/10.1007/s40747-020-00262-0} (arXiv:1906.06660).

\end{thebibliography}

\end{document}

